\id{МРНТИ 52.13.17, 52.13.15, 52.01.75}{}
\vspace{1em}
\begin{header}
\swa{Д.К.~Бекбергенов, И.Н.~Савич, А.А.~Зейнуллин, Ж.~Бимурат, Р.К.~Жанакова}{ОПТИМИЗАЦИЯ КАЧЕСТВЕННЫХ ПАРАМЕТРОВ ДОБЫЧИ РУД В ЗОНАХ ОБРУШЕНИЯ ЖЕЗКАЗГАНСКОГО МЕСТОРОЖДЕНИЯ НА ОСНОВЕ ЭММ}

\tsp{1}Д.К.~Бекбергенов\envelope,
\tsp{2}И.Н.~Савич,
\tsp{3}А.А.~Зейнуллин,
\tsp{1}Ж.~Бимурат,
\tsp{4}Р.К.~Жанакова\envelope
\end{header}

\begin{affil}
\tsp{1}Институт горного дела им. Д. А. Кунаева, Алматы, Казахстан,

\tsp{2}Горный институт НИТУ МИСиС, Москва, Российская Федерация,

\tsp{3}АО «Казахский университет технологии и бизнеса им.К.Кулажанова», Астана, Казахстан,

\tsp{4}Казахский национальный исследовательский технический университет им. К. И. Сатпаева, Алматы, Казахстан

\corrauthor{Корреспондент-автор: kdbekbergen@mail.ru, zhanakova\_raisa@mail.ru}
\end{affil}

Статья посвящена решению стратегической задачи вовлечения в отработку
нетронутых залежей шахты «Анненская» Восточно-Жезказганского рудника,
находящихся в зонах обрушения с активным сдвижением массива. В условиях,
где традиционные геотехнологии неэффективны, обосновано применение
сплошной системы разработки с обрушением и торцевым выпуском руды,
усиленной комбинированной крепью СВП-22.

Научная новизна исследования заключается в разработке и применении
экономико-математи\-ческой модели (ЭММ) критерием оптимизации которой
является получение максимальной прибыли на основе регрессионного анализа
с регуляризацией.

Использование стандартизированной формы уравнения прибыли позволило
авторам впервые корректно сопоставить разнородные факторы (геометрию
залежи, объемы добычи и качественные показатели). Модель выявила
нелинейный характер влияния потерь и разубоживания на финансовый
результат, а также определила роль тектонического сдвига как
модификатора чувствительности прибыли к изменению угла залегания рудных
тел. В результате исследования установлены критические зависимости:
получено семейство вогнутых кривых (поверхностей отклика), описывающих
экстремумы прибыли. Доказано, что при увеличении угла падения свыше 30
градуса область рентабельности резко сужается, а поддержание
сверхнизкого уровня потерь становится экономически нецелесообразным
из-за опережающего роста эксплуатационных затрат. Установлен
технологический предел применения метода - угол падения до 45 градуса.
Практическая значимость работы состоит в создании компактного
математического инструмента, позволяющего определять точки
безубыточности и выбирать оптимальные стратегии размещения
буро-доставочных штреков в условиях высокой неопределенности нарушенного
массива.

{\bfseries Ключевые слова}: экономико-математическая модель, нетронутые
залежи, математическое моделирование, система разработки, прибыль, зона
обрушения, буро-доставочный штрек.
\vspace{1em}
\begin{header}
ЭММ НЕГІЗІНДЕ ЖЕЗҚАЗҒАН КЕН ОРНЫНЫҢ ОПЫРЫЛУ АЙМАҚТАРЫН ИГЕРУ КЕЗІНДЕГІ ПАЙДАНЫ БОЛЖАУ

\tsp{1}Д.К.~Бекбергенов\envelope,
\tsp{2}И.Н.~Савич,
\tsp{3}А.А.~Зейнуллин,
\tsp{1}Ж.~Бимурат,
\tsp{4}Р.К.~Жанакова\envelope
\end{header}

\begin{affil}
\tsp{1}А. Қонаев атындағы тау-кен ісі институты, Алматы, Қазақстан,

\tsp{2}НИТУ МИСИС, Москва, Россия,

\tsp{3}«Қ. Құлажанов атындағы Қазақ технология және бизнес университеті» АҚ, Астана. Қазақстан,

\tsp{4}Қ.И. Сәтбаев атындағы Қазақ ұлттық техникалық зерттеу университеті, Алматы, Қазақстан,

е-mail: kdbekbergen@mail.ru, zhanakova\_raisa@mail.ru
\end{affil}

Мақала массивтің белсенді жылжу аймақтарында орналасқан «Анненская»
шахтасының тиілмеген кен жатындарын игеруге тартудың стратегиялық
міндетін шешуге арналған. Дәстүрлі геотехнологиялар тиімсіз болатын
жағдайларда, СВП-22 құрамдастырылған бекітпесімен күшейтілген, кенді
опыру және шеткі шығарумен жаппай өңдеу жүйесін қолдану негізделген.
Зерттеудің ғылыми жаңалығы регуляризацияланған регрессиялық талдау
негізінде экономикалық-математикалық модельді (ЭММ) әзірлеу және қолдану
болып табылады.

Пайда теңдеуінің стандартталған формасын пайдалану авторларға алғаш рет
әртүрлі факторларды (жатынның геометриясын, өндіру көлемін және сапалық
көрсеткіштерді) дұрыс салыстыруға мүмкіндік берді. Модель кеннің ысырабы
мен кедейленуінің қаржылық нәтижеге әсерінің сызықты емес сипатын
анықтады, сондай-ақ тектоникалық ығысудың кен денелерінің жату бұрышының
өзгеруіне пайданың сезімталдық модификаторы ретіндегі рөлін айқындады.

Зерттеу нәтижесінде критикалық тәуелділіктер анықталды: пайданың
экстремумдарын сипаттайтын ойыс қисықтар (жауап беру беттері) жиынтығы
алынды. Кеннің құлау бұрышы 30 градустан асқанда рентабельділік аймағы
күрт тарылатыны, ал ысыраптың өте төмен деңгейін сақтау пайдалану
шығындарының жедел өсуіне байланысты экономикалық тұрғыдан тиімсіз
болатыны дәлелденді. Әдісті қолданудың технологиялық шегі құлау бұрышы
45 градусқа дейін екені белгіленді. Жұмыстың практикалық маңыздылығы
бұзылған массивтің жоғары белгісіздік жағдайында шығынсыздық нүктелерін
анықтауға және бұрғылау-жеткізу штректерін (БЖШ) орналастырудың оңтайлы
стратегияларын таңдауға мүмкіндік беретін ықшам математикалық құралды
жасаудан тұрады.

{\bfseries Түйін сөздер:} экономикалық-математикалық модель, тиілмеген кен
оырндары, математикалық модельдеу, өңдеу жүйесі, пайда, тасымалдау
штрегі
\vspace{1em}
\begin{header}
PROFIT FORECASTING FOR MINING DEPOSITS IN CAVING ZONES OF THE ZHEZKAZGAN FIELD BASED ON EMM

\tsp{1}D.K.~Bekbergenov\envelope,
\tsp{2}I.N.~Savich,
\tsp{3}А.А.~Zeinullin,
\tsp{1}Zh.~Bimurat,
\tsp{4}R.K.~Zhanakova\envelope
\end{header}

\begin{affil}
\tsp{1}D.A. Kunaev Institute of Mining, Almaty, Kazakhstan,

\tsp{2}Mining Institute of NUST MISIS, Moscow, Russian Federation, Moscow, Russia,

\tsp{3}«K. Kulazhanov Kazakh University of Technology and Business» JSC, Astana, Kazakhstan,

\tsp{4}K.I. Satpayev Kazakh National Research Technical University, Almaty, Kazakhstan,

е-mail: kdbekbergen@mail.ru, zhanakova\_raisa@mail.ru
\end{affil}

The article addresses the strategic challenge of extracting untouched
deposits at the Annenskaya mine (East Zhezkazgan Mine), which are
located within caving zones characterized by active rock mass
dis\-placement Given the inefficiency of traditional geotechnologies, the
use of a continuous mining system with caving and sublevel (end)
drawing, reinforced with SVP-22 combined support, is justified.

The scientific novelty of the research lies in the development and
application of an economic-mathema\-tical model (EMM) based on regression
analysis with Lasso regularization.

The use of a standardized profit equation allowed the authors for the
first time to correctly compare heterogeneous factors (deposit geometry,
production volumes, and quality indicators). The model revealed the
non-linear nature of the impact of losses and dilution on the financial
result and identified the role of tectonic shift as a modifier of profit
sensitivity to changes in the dip angle. As a result of the study,
critical dependencies were established: a family of concave curves
(response surfaces) describing profit extrema was obtained. It is proved
that with an increase in the dip angle beyond 30 degrees, the
profitability zone narrows sharply, and maintaining an ultra-low level
of losses becomes economically inexpedient due to the outstripping
growth of operating costs. The technological limit of the method is set
at a dip angle of up to 45 degrees. The practical significance of the
work lies in the creation of a compact mathematical tool that allows
determining break-even points and choosing optimal strategies for
locating drilling-and-delivery drifts (DDD) under high uncertainty of a
disturbed rock mass.

{\bfseries Keywords:} economic-mathematical model, intact ore bodies,
mathematical modeling, mining system, profit, drilling and transport
tunnels

\begin{multicols}{2}
{\bfseries Введение.} Современный этап эксплуатации Жезказганского
месторождения характеризуется переходом к освоению запасов в
экстремально сложных геомеханических условиях. Одной из наиболее острых
проблем является наличие значительных объемов «нетронутых» залежей,
оказавшихся заблокированными в зонах обрушения и мульдах сдвижения
массива, в частности на шахте «Анненская» ВЖР. В таких условиях
традиционные методы камерно-столбовой отработки становятся
неэффективными и опасными из-за высокого риска неконтролируемых
обрушений, неопределенности морфологии рудных тел и неизбежно высоких
потерь полезного компонента. Вовлечение этих трудноизвлекаемых ресурсов
в рентабельную эксплуатацию требует внедрения принципиально новых
технологических решений и инструментов точного экономического
прогнозирования.

В качестве адекватного ответа на сложную геомеханическую обстановку
{[}1{]} предлагается переход к технологии сплошной системы разработки с
обрушением и торцевым выпуском руды. Данное решение позволяет исключить
нахождение персонала в опасных зонах за счет использования
буро-доставочных штреков (БДШ) малого сечения. Для обеспечения
устойчивости выработок в нарушенном массиве применяется комбинированная
крепь, сочетающая анкерное крепление и металлический спецпрофиль СВП-22
{[}2{]}. Однако принятие решений об отработке конкретных блоков
осложняется многофакторностью процесса: на финальную прибыль влияют не
только горно-геологические параметры (мощность, угол падения), но и
переменные показатели разубоживания, потерь и растущие затраты на
поддержание выработок.

Параллельно с технологическим перевооружением, развитие вычислительной
техники и внедрение геоинформационных систем (ГИС) открывают возможности
для высокоточного проектирования. Накопленный массив геологических
данных по Анненскому рудному полю Жезказганского месторождения послужил
фундаментом для создания экономико-математической модели (ЭММ),
адекватно отражающей пространственную геометрию рудных тел даже в
условиях зон обрушения. Использование многофакторных квадратичных
моделей позволяет установить четкую математическую зависимость между
параметрами отработки и финансовым результатом, обеспечивая научное
обоснование выбора наиболее эффективной схемы размещения БДШ.

Целью данной работы является выбор и обоснование оптимальных
качественных показателей добычи руды на основе прогнозной модели
получения максимальной расчетной прибыли на ЭММ для минимизации
геологических рисков и максимизации извлечения запасов. Благодаря
построению уточненных блочных моделей, авторами реализована возможность
не только прецизионного подсчета объемов руды при сложных
горно-геологических условиях, но и проведения сценарного анализа
рентабельности. Такой комплексный подход сочетание безопасной технологии
торцевого выпуска и математической оптимизации параметров добычи
позволяет эффективно вовлекать в эксплуатацию ранее оставленных
нетронутых залежей, обеспечивая устойчивое развитие Жезказганского
месторождения на его завершающей стадии {[}3{]}.

Для эффективного извлечения запасов, сосредоточенных в нетронутых
залежах Анненского рудного поля, обосновано применение сплошной системы
разработки с обрушением и торцевым выпуском руды {[}4{]}. Данная
технология наиболее адаптивна к условиям нарушенного массива, так как
позволяет вести добычу без формирования открытого очистного
пространства. Основным конструктивным параметром блока здесь выступает
мощность рудной залежи, определяющая его высоту, в то время как шаг
обрушения слоя принимается равным величине линии наименьшего
сопротивления (ЛНС) взрываемого веера скважин. Для сохранения
устойчивости выработок и минимизации сейсмического воздействия на
вмещающий массив, отбойка производится путем взрывания одного веера с
применением схем короткозамедленного детонирования.

Подготовка выемочных участков реализуется через проходку
буро-доставочных штреков (БДШ) для каждой панели блока. Выбор способа их
поддержания напрямую зависит от локальных горно-геологических условий:
выработки проводятся либо без крепления, либо с применением
комбинированной системы усиления. В последнем случае используется
анкерное крепление кровли в сочетании со стальными арочными рамами из
специального взаимозаменяемого профиля (СВП-22) {[}5-7{]}. Схема
бурения, представленная на разрезе Б-Б (рисунок 1), предусматривает
веерное расположение скважин по рудному телу с обязательным перебуром до
двух метров во вмещающие породы, что гарантирует полноту извлечения и
надежную фиксацию анкеров в устойчивой части массива.
\end{multicols}

\fig{g/image63}[Рис.1 - разрез Б - Б при проходке БДШ по нетронутой рудной
залежи рекомендуемой сплошной системой обрушения и с торцевым выпуском
рудной массы в условиях зоны обрушения с размещением схемы бурения
вееров скважины по руде и с перебуром до 2 метров во вмещающем породном
массиве и анкерного крепления кровли БДШ]

\begin{multicols}{2}
Особое внимание в работе уделено применению металлических трехзвенных
податливых крепей из профиля СВП-17 и выше {[}6,7{]}. Конструктивно
такая крепь состоит из отдельных арок (включающих верхняк и две стойки),
устанавливаемых с шагом 0,5--1,2 м и объединенных в единый каркас
межрамными стяжками. Использование желобчатого спецпрофиля и узлов
податливости на основе скоб с планками позволяет крепи воспринимать
значительные деформации массива без потери несущей способности. Для
предотвращения вывалов между арками устанавливаются деревянные,
железобетонные или металлические затяжки.

Внедрение данных конструкций рекомендуется для выработок, проводимых в
слабых, сильнотрещиноватых и неустойчивых породах с коэффициентом
крепости f = 3-9 при углах наклона до 30°. Выбор конкретного типоразмера
профиля (от СВП-14 до СВП-33{[}2{]}) осуществляется исходя из расчетной
величины ожидаемых смещений массива в зоне мульдыижения шахты
«Анненская». Такой комплексный подход к креплению и подготовке БДШ
обеспечивает безопасность персонала и оборудования при последующем
выполнении буровзрывных работ и погрузке отбитой горной массы
самоходными ПДМ.

Подготовительный цикл работ при реализации предложенной системы
предполагает проведение сети буро-доставочных штреков по простиранию
рудного тела. Проходка БДШ осуществляется от существующих или специально
возведенных разрезных штреков (ортов) и продолжается до полной границы
выклинивания залежи или на расчетную длину панели. С целью обеспечения
планомерного извлечения запасов штреки закладываются с интервалом 12
метров, начиная от верхней границы блока с последовательным развитием
фронта работ по падению залежи.

Учитывая геологическую изменчивость Жезказганского месторождения,
технологическая схема адаптирована под различные геомеханические
сценарии. Основными переменными факторами выступают мощность рудной
залежи (в диапазоне 6, 8, 10 и 12 м) и углы её падения, составляющие 30°
и 50°. Исследованием охвачены типовые варианты отработки, учитывающие
эти характеристики: например, наклонные залежи мощностью 6 м при угле
падения 30°, а также более крутые пласты мощностью 9 м с углом залегания
50° . Данная система разработки рекомендуется в вариантах по мощности
залежи (6, 8, 10, 12 м) и углах падениях залежи 30 и 50 градусов (см.
рисунки 2 и 3).
\end{multicols}

\begin{figs}
\fig[0.51\textwidth][5cm]{g/image64}[Рис.2 - Вариант технологической схемы отработки
нетронутых залежей при мощности 6 м и угле падения рудного тела - 30º]
\fig[0.39\textwidth][5cm]{g/image65}[Рис.3 - Вариант технологической схемы отработки нетронутых
залежей при мощности 6 м и угле падения рудного тела - 50º]
\end{figs}

\begin{multicols}{2}
Представленные графические модели (рис.2 и 3) иллюстрируют рациональное
расположение БДШ и конфигурацию очистной выемки, обеспечивающие
устойчивость массива в условиях зоны обрушения. Стоит отметить, что
данные варианты являются базовыми: в ходе дальнейших исследований
планируется расширение вариативности технологических схем для охвата
всего спектра горно-геологических условий шахты «Анненская». Это
позволит сформировать универсальный алгоритм выбора параметров системы
разработки, исходя из конкретных метрических данных каждого блока.

{\bfseries Материалы и методы.} \emph{Экономико-математическое
моделирование и оптимизация качественных параметров добычи критерием
максимальной прибыли Жезказганского месторождения на основе
многофакторных квадратичных моделей.}

В процессе выполнения основных этапов математического моделирования были
рассмотрены основные факторы, влияющие на точность получаемой модели и
приведены два основных фактора -- потери и разубоживание, которые
определяют необходимую и достаточную точность модели месторождения. Учет
потерь, как неизвлеченные из недр запасы и разубоживания, как засорение
руды в добыче полезных ископаемых контролирует то, как полно извлекается
ценное сырье из недр и как снижается его качество из-за примесей пустой
породы, что критично для оценки запасов, планирования и экономической
эффективности {[}8,9{]}.

На данном исследовании рассматривается квадратичная факторная модель
прибыли, в основу которой положен мультипликативный принцип: итоговая
прибыль определяется как произведение базовой прибыли, основанную на
принципе:

Прибыль = Базовая прибыль х Влияние угла наклона залежи х Влияние потерь
и разубоживания, где каждый множитель имеет квадратичную форму с
максимумом в целевой точке.

Оптимизирующим критерием выбора лучших качественных параметров добычи
руд является получение максимальной прибыли: П = Ц -- С → max,
где, Ц -- ценность руды.

В показатель себестоимости добычи С суммируется ущерб от потерь и
разубоживания руды.

Общая формула прибыли:

\begin{equation}
П(\theta,L,D) = П_{ref} \times \Phi_{\theta}(\theta) \times \Psi_{LD}(L,D)
\end{equation}

Влияние угла падения (\(\Phi_{\theta}\)):

\begin{equation}
\Phi_{\theta} = \max\left( 0;1 - k_{\theta} \times \left( \frac{\theta - \theta^{*}}{\Delta \theta} \right)^{2} \right)
\end{equation}

Нормализация отклонения
\(\frac{\theta - \theta^{*}}{\Delta \theta}\) приводит
отклонение к безразмерному виду.

Квадратичный штраф
\(\left( \frac{\theta - \theta^{*}}{\Delta \theta} \right)^{2}\)отражает,
что убытки растут быстрее линейного.

Коэффициент жесткости \(k_{\theta}\) определяет крутизну падения (0--1).

Обрезка нулем \({max}(0;\ldots)\ \)исключает отрицательные значения.

Влияние потерь и разубоживания (\(\Psi_{LD}(L,D)\))
\end{multicols}

\begin{equation}
\Psi_{LD}(L,D) = {max}\left( 0;1 - a \times \left( \frac{L - L^{*}}{\Delta L} \right)^{2} - b \times \left( \frac{D - D^{*}}{\Delta D} \right)^{2} - c \times \left( \frac{L - D}{\Delta \Delta } \right)^{2} \right)
\end{equation}

\begin{multicols}{2}
Теория построения трех слагаемых:

- ущерб от потерь:
\(a \times \left( \frac{L - L^{*}}{\Delta L} \right)^{2}\), где
\(L\) - фактические потери (доля), \(L^{*}\) - целевой уровень потерь,
\(\Delta L\) - допустимое отклонение, \(a\) - вес этого фактора
(0.25).

- ущерб от разубоживания:
\(b \times \left( \frac{D - D^{*}}{\Delta D} \right)^{2}\), где
\(D\) - фактическое разубоживание (доля), \(D^{*}\) - целевой уровень
разубоживания, \(\Delta D\) -- допустимое отклонение, \(b\) -
вес (0.30 больше, чем у \(a\) - разубоживание критичнее).

- штраф за разбаланс между потерями и разубоживанием:
\(c \times \left( \frac{L - D}{\Delta \Delta } \right)^{2}\),
где \((L - D)\) -- разность между потерями и разубоживанием,
\(\Delta \Delta \) - допустимый разбаланс, \(c\) -- вес
(0.15).

Физический смысл: сильная диспропорция между потерями и разубоживанием
указывает на системные проблемы технологии. Математическая постановка
точка безубыточности \(П(\theta,L,D) = 0\). Условия: либо
\(\Phi_{\theta} = 0\), либо \(\Psi_{LD}(L,D) = 0\).

Решение для угла

\begin{equation}
1 - k_{\theta} \times \left( \frac{\theta - \theta^{*}}{\Delta \theta} \right)^{2} = 0
\end{equation}

\begin{equation}
\left( \frac{\theta - \theta^{*}}{\Delta \theta} \right)^{2} = \frac{1}{k_{\theta}}
\end{equation}

\begin{equation}
\theta = \theta^{*} \pm \frac{\Delta \theta}{\sqrt{k_{\theta}}}
\end{equation}

Решение для потерь/разубоживания. Уравнение \(\Psi_{LD}(L,D) = 0\)
квадратное относительно \(L\) и \(D\). Для \(L\) при фиксированном \(D\)
\end{multicols}

\begin{equation}
1 - a \times \left( \frac{L - L^{*}}{\Delta L} \right)^{2} - b \times \left( \frac{D - D^{*}}{\Delta D} \right)^{2} - c \times \left( \frac{L - D}{\Delta \Delta } \right)^{2} = 0
\end{equation}

приводится к виду \(A \times L^{2} + B \times L + C = 0\), где

\begin{equation}
A = \frac{a}{{\Delta L}^{2}} + \frac{c}{{\Delta \Delta }^{2}}
\end{equation}

\begin{equation}
B = - \frac{2a \times L^{*}}{{\Delta L}^{2}} - \frac{2c \times D}{{\Delta \Delta }^{2}}
\end{equation}

\begin{equation}
C = 1 - a \times \left( \frac{L^{*}}{\Delta L} \right)^{2} - b \times \left( \frac{D - D^{*}}{\Delta D} \right)^{2} - c \times \left( \frac{D}{\Delta \Delta } \right)^{2}
\end{equation}

\begin{multicols}{2}
Ключевым допущением модели является оптимальность данной точки, что
обуславливает равенство нулю первых производных (градиента) и позволяет
сфокусироваться на матрице Гессе для описания квадратичного штрафа за
отклонение от целевых показателей.

\emph{Экономико-математическое моделирование прибыли на основе
регрессионного анализа с Lasso-регуляризацией.}

Далее исследования посвящен разработке регрессионной модели прибыли,
которая расширяет базовые представления о зависимостях за счет включения
качественных характеристик руды и использования современных методов
машинного обучения. В отличие от классических подходов, здесь
формируется расширенная полиномиальная модель второго порядка,
интегрирующая не только угол падения (\(\theta)\) , потери (L) и
разубоживание (D), но и такие специфические для Жезказганского
месторождения параметры, как балансовые запасы (Q), содержание меди (M)
и коэффициент извлечения (\(\xi\)). Для обеспечения математической
устойчивости модели в условиях ограниченной выборки применяется
Lasso-регрессия (L1-регуляризация), которая автоматически отсеивает
малозначимые факторы и устраняет мультиколлинеарность между
полиномиальными членами {[}10-12{]}.

На первом этапе формируется полиномиальная регрессионная модель второго
порядка, описывающая зависимость прибыли от геолого-технологических
параметров отработки:
\end{multicols}
\vspace{-1cm}
\begin{multline}
П = \beta_{0} + \beta_{1}\theta + \beta_{2} L + \beta_{3} D + \beta_{4} S + \beta_{5}\theta^{2} + \beta_{6} L^{2} + \beta_{7} D^{2} + \beta_{8}\theta L \\
+ \beta_{9}\theta D + \beta_{10} LD + \beta_{11}\theta S + \beta_{12} LS + \beta_{13} DS + \epsilon
\end{multline}

\begin{multicols}{2}
где, \(П\) - прибыль, USD (\$);

\(\theta\) - угол падения залежи, градусы;

\(L\) - потери полезного ископаемого (доли);

\(D\) - разубоживание (доли);

\(S\) - сдвиг (0 - отсутствует, 1 - присутствует);

\(\beta_{i}\) - параметры регрессионной модели;

\(\epsilon\) - случайная ошибка (ошибка модели).

Для Жезказганского меднорудного месторождения прибыль формируется не
только геометрией и потерями, но и качественными характеристикам руды и
запасов. В связи с этим модель расширяется следующими параметрами:

\(Q_{bal}\) - балансовые запасы, т;

\(c_{Cu}\) - среднее содержание меди, доли или \%;

\(\xi\) - коэффициент сквозного извлечения.

Расширенная регрессионная модель принимает вид:
\end{multicols}
\vspace{-1cm}
\begin{multline}
П = \beta_{0} + \beta_{1}\theta + \beta_{2}L + \beta_{3}D + \beta_{4}S + \beta_{5}Q_{bal} + \beta_{6}c_{Cu} + \beta_{7}\xi \\
+ \beta_{8}\theta^{2} + \beta_{9}L^{2} + \beta_{10}D^{2} + \beta_{11}\theta L + \beta_{12}\theta D + \beta_{13}LD + \beta_{14}\theta S + \epsilon
\end{multline}

\begin{multicols}{2}
Для обеспечения устойчивости оценок и корректной интерпретации
коэффициентов выполняются следующие операции:

1. Перевод процентных показателей L и D в доли.

2. Стандартизация всех числовых признаков:
\(z_{i} = \frac{x_{i} - \mu_{i}}{\sigma_{i}}\).

3. Разделение выборки на обучающую (80\%) и тестовую (20\%).

Для оценки коэффициентов применяется Lasso-регрессия, минимизирующая
функционал:

\begin{equation}
\frac{1}{n}\sum_{i = 1}^{n}\left( П_{i} - {\widehat{П}}_{i} \right)^{2} + \lambda\sum_{j}^{}\left| \beta_{j} \right|
\end{equation}

Использование Lasso позволяет: автоматически отбирать значимые линейные,
квадратичные и перекрёстные члены; устранить мультиколлинеарность,
возникающую при полиномиальном расширении; получить компактную и
интерпретируемую модель прибыли.

После стандартизации признаков и применения Lasso-регрессии для
Жезказганского месторождения была получена следующая модель прибыли:
\end{multicols}
\vspace{-0.7cm}
\begin{equation}
\widetilde{П} = \beta_{0} - 0.62\widetilde{L} - 0.48\widetilde{D} + 0.21\widetilde{\theta} - 0.31{\widetilde{L}}^{2} - 0.27{\widetilde{D}}^{2} + 0.18{\widetilde{\theta}}^{2} - 0.35\widetilde{L}\widetilde{D} + 0.16\widetilde{\theta} - 0.12\widetilde{\theta}S
\end{equation}

\begin{multicols}{2}
где, тильда обозначает стандартизированные переменные. Это
адаптированная модель именно для Жезказганского месторождения.

{\bfseries Результаты и обсуждение.} Анализ диапазонов входных данных
позволяет определить операционное окно рентабельности и установить
границы чувствительности модели к внешним и технологическим факторам. В
контексте геомеханических условий Жезказганского месторождения
исследуемый диапазон углов падения охватывает значения от
30\tsp{0} до 51\tsp{0}, при этом наиболее
благоприятным и экономически эффективным признан интервал от 300 до
45\tsp{0}. Для технологических параметров установлены
вариативные границы: потери руды (L) варьируются в пределах 19--30\%,
что определяет допустимое отклонение D-L на уровне 0.06 (или половину
амплитуды диапазона), в то время как разубоживание (D) демонстрирует
более широкую область значений от 23\% до 38\% с расчетным допуском D =
0.08. Подобная дифференциация допусков математически обосновывает
различную степень влияния каждого фактора на итоговую прибыль,
подтверждая, что стабильность системы в большей степени зависит от
контроля показателей разубоживания.

Более высокий вес разубоживания (0.30 против 0.25 у потерь)
математически подтверждает, что для условий Жезказгана засорение руды
пустой породой является более критичным фактором, чем недовыемка
запасов. Для адаптации модели к условиям Жезказганского месторождения
использованы эмпирические данные, согласно которым базовая прибыль
составляет 1.15х10\tsp{6}, а целевые уровни потерь и
разубоживания установлены на отметках 22\% и 27\% соответственно. Анализ
модели позволяет аналитически определять точки безубыточности через
решение квадратных уравнений относительно каждого фактора.

Предложенная методология аппроксимации 2-го порядка позволяет не только
визуализировать «поверхности отклика» прибыли в 3D-пространстве, но и
количественно оценить системные риски технологии через штраф за
разбаланс (L−D). Это обеспечивает инженерно-технический персонал
инструментом для оперативной корректировки планов горных работ с целью
удержания параметров в «куполе рентабельности», ограниченном допустимыми
отклонениями по потерям (±6\%) и разубоживанию (±8\%).

На основе разработанной математической модели (рисунок 4) и
установленных калибровочных параметров, представленный рисунок
визуализирует комплексную зависимость прибыли Жезказганского
месторождения от ключевых факторов отработки. Центральное место в
анализе занимает базовая прибыль, которая при оптимальных условиях
достигает максимума в размере около 1.15х10\tsp{6}. Графики
наглядно демонстрируют квадратичный характер функции прибыли: любое
отклонение от целевых значений угла падения (\(\theta\) =
30\tsp{0}), потерь (L = 22\%) или разубоживания (D = 27\%)
ведет к параболическому снижению финансового результата.
\end{multicols}

\fig[0.75\textwidth]{g/image66}[Рис.4 - Результаты многофакторного квадратичного моделирования
прибыли для условий Жезказганского месторождения]
\fig[0.7\textwidth]{g/image67}[Рис.5 - Результаты математического моделирования в ззависимости
прибыли от уровня потерь рудного тела с фиксацией точек безубыточности]

\begin{multicols}{2}
Особое внимание уделено влиянию геометрии залежи и критическим порогам
эксплуатации. Моделирование показывает, что при достижении критического
угла падения в 60.0° прибыль полностью нивелируется квадратичным
штрафом, независимо от уровня технологических потерь. Аналогично,
графики «Прибыль vs Разубоживание» фиксируют точку безубыточности на
отметке 31.3\%, подтверждая заложенный в модель более высокий вес штрафа
за разубоживание (\(\Psi_{D}\)= 0.30) по сравнению с потерями
(\(\Psi_{L}\)= 0.25). Это подчеркивает, что для экономики Жезказгана
чистота добываемой руды является более значимым фактором, чем полнота
извлечения запасов.

Анализ системной устойчивости технологии отражен на графике зависимости
прибыли от разбаланса (L-D). В сочетании с представленными поверхностями
отклика, которые при увеличении угла до
50\tsp{0}--51\tsp{0} практически полностью
уходят в область отрицательных значений, поэтому максимальный угол
падения рекомендуется до 45\tsp{0} в заданных
горно-геологических условиях (рис 5--6).
\end{multicols}

\begin{figs}
    \fig[0.4\textwidth]{g/image68}
    \fig[0.4\textwidth]{g/image69}
\end{figs}
\vspace{0.5cm}
\begin{figs}[Рис.6 - 3D экономическая-математическая модель области
рентабельности отработки при угле наклона залежи]
    \fig[0.4\textwidth]{g/image70}
    \fig[0.4\textwidth]{g/image71}
\end{figs}

\begin{multicols}{2}
{\bfseries Выводы.} Научная новизна данного метода заключается в
использовании стандартизированной формы уравнения прибыли, что позволяет
корректно сопоставлять влияние разнородных факторов (градусов, тонн и
процентов). Полученная модель фиксирует, что наибольшее деструктивное
воздействие на финансовый результат оказывают потери и разубоживание,
причем их влияние носит ярко выраженный нелинейный характер,
усиливающийся при определенных комбинациях геометрических параметров.
Фактор тектонического сдвига модели выступает как модификатор
чувствительности, изменяющий реакцию прибыли на изменение угла
залегания, что полностью соответствует геологическим реалиям
Жезказганского региона.

Интерпретация результатов регрессионного моделирования подтверждает
выводы визуального анализа поверхностей отклика: при малых значениях
потерь влияние угла падения проявляется умеренно, однако по мере роста
технологических отклонений чувствительность системы к геометрии залежи
резко возрастает. Использование регуляризации позволило получить
компактную и интерпретируемую формулу, которая служит надежным
инструментом для проведения анализа чувствительности и поиска
оптимальных стратегий отработки месторождения в условиях высокой
неопределенности.

Представленные графические зависимости иллюстрируют нелинейный характер
изменения прибыли в функции от коэффициента потерь рудного тела при
варьировании угла падения залежи. Математические модели описываются
семейством вогнутых кривых (парабол), что свидетельствует о наличии
экстремума оптимального уровня технологических потерь, при котором
достигается максимальный финансовый результат. С ростом угла наблюдается
выраженная деградация доходности: кривые смещаются в область
отрицательных значений, а амплитуда максимальной прибыли существенно
сокращается, что объясняется усложнением условий отработки и ростом
удельных затрат.

Расчеты показали, что при увеличении ценности руды оптимум прибыли
смещается в сторону снижения ущерба от потерь и разубоживания руды.

Маркированные точки на графике определяют критические значения
коэффициента потерь, разделяющие области рентабельной и убыточной
добычи. Для пологих залежей точка безубыточности достигается при
минимальных значениях потерь (L в приделах 0,21), обеспечивая широкий
диапазон эффективной эксплуатации. Однако при увеличении угла наклона от
50\tsp{0} область рентабельности сужается до критического
минимума, смещаясь в сторону более высоких значений потерь. Это
указывает на то, что для крутопадающих тел поддержание низкого уровня
потерь становится экономически нецелесообразным из-за опережающего роста
эксплуатационных расходов. По данному методу также максимальный угол
падения рекомендуется до 45\tsp{0} в заданных
горно-геологических условиях.
\end{multicols}

\begin{center}
{\bfseries Литература}
\end{center}

\begin{refs}
1. Гельманова З. С., Базаров Б. А., Мезенцева А. В., Конакбаева А. Н.,
Толешов А.К. Управление производственной безопасностью на
горнодобывающих предприятиях Казахстана. Караганда: Изд-во КарГТУ, 2017.
124 с.

2. Каспарьян Э.В., Кузнецов А. А. Геомеханика: учебник для вузвузов.
-Москва: Юрайт- 2023. - 245 с. ISBN 978-5-86185-901-1

3. Бекбергенов Д.К., Савич И. Н., Зейнуллин А.А. Технологическое решение
для подземного освоения нетронутых залежей, оставленных в зоне обрушения
Жезказганского месторождения // Горный информационно-аналитический
бюллетень. - 2025. - № 6 (спецвыпуск 10). -- 20 с.

4. Крыкпышев Б. Проблемы и пути инновационного развития горной
промышленности: презентационный доклад // IV международная
научно-практическая конференция «Геотехника-2013». - Алматы, 2013.

5. Бекбергенов Д. К. Итоговый отчет по теме НИР «Исследование
соответствия определения параметров и системы разработки в условиях
шахты „Анненская`` Восточно-Жезказганского рудника». -- Алматы: ТОО
«КазНИИЦветмет», 2022. - 347 с.

6. Бекбергенов Д.К., Зейнуллин А.А. Развитие технологии повторной добычи
руды из междукамерных целиков, оставленных в зонах обрушения на
завершающей стадии отработки Жезказганского месторождения // Горный
информационно-аналитический бюллетень. -- 2025. -- № 2 (спецвыпуск 4).
-- 20 с.

7. Demin V., Kalinin A., Tomilova N., Tomilov A., Akpanbayeva A.,
Shokarev D., Popov A., Advanced digital modeling of stress--strain
behavior in rock masses to ensure stability of underground mine
workings//Civil Engineering Journal. - 2025.-Vol.11(3). -P.1072-1087.
\href{https://doi.org/10.28991/CEJ-2025-011-03-014}{DOI
10.28991/CEJ-2025-011-03-014}.

8. Аханов Т. М. Создание высокоэффективных технологий разработки
месторождений в условиях подземных рудников ТОО «Корпорация Казахмыс» //
Горный журнал Казахстана. - 2009. - № 11. - С.26--30.

9. Герасименко В. И., Макаров А. Б., Зотеев О. В. Анализ
геомеханического состояния выработанных пространств Жезказганского
месторождения. Часть 1. - Караганда - Жезказган: ГГУ ТОО «Корпорация
Казахмыс», 2010. - 52 с.

10. Низамова Г. З., Гайфуллина М. М. Корреляционно-регрессионный анализ
эффективности использования инвестиционных ресурсов нефтяной компании //
Вестник УГНТУ. Наука, образование, экономика. Серия: Экономические
науки. - 2021. - Т.9(3).-C.27-38 DOI
\href{https://doi.org/10.26425/2309-3633-2021-9-3-27-38}{10.26425/2309-3633-2021-9-3-27-38}.

11. Самсонов А. В., Васильев И. В. Анализ инвестиционной
привлекательности предприятий нефтегазовой отрасли в России //
Московский экономический журнал. -- 2020. - № 2. - С.439--450. DOI
10.24411/2413-046x-2020-10088.

12. Hastie T., Tibshirani R., Friedman J. The Elements of Statistical
Learning: Data Mining, Inference, and Prediction. - 2nd ed. - New York:
Springer. - 2009. - 745 p.
\end{refs}

\begin{center}
{\bfseries References}
\end{center}

\begin{refs}
1. Gel' manova Z. S., Bazarov B. A., Mezenceva A. V.,
Konakbaeva A. N., Toleshov A.K. Upravlenie proizvodstvennoj
bezopasnost' ju na gornodobyvajushhih predprijatijah
Kazahstana. Karaganda: Izd-vo KarGTU, 2017.124 s. {[}in Russian{]}

2. Kaspar' jan Je.V., Kuznecov A. A. Geomehanika:
uchebnik dlja vuzov. - Moskva: Jurajt- 2023. - 245 s. ISBN
978-5-86185-901-1. {[}in Russian{]}

3. Bekbergenov D.K., Savich I. N., Zejnullin A.A. Tehnologicheskoe
reshenie dlja podzemnogo osvoenija netronutyh zalezhej, ostavlennyh v
zone obrushenija Zhezkazganskogo mestorozhdenija // Gornyj
informacionno-analiticheskij bjulleten'. - 2025. - № 6
(specvypusk 10). - 20 s. {[}in Russian{]}

4. Krykpyshev B. Problemy i puti innovacionnogo razvitija gornoj
promyshlennosti: prezentacionnyj doklad // IV mezhdunarodnaja
nauchno-prakticheskaja konferencija «Geotehnika-2013». - Almaty, 2013.
{[}in Russian{]}

5. Bekbergenov D. K. Itogovyj otchet po teme NIR «Issledovanie
sootvetstvija opredelenija parametrov i sistemy razrabotki v uslovijah
shahty „Annenskaja`` Vostochno-Zhezkazganskogo rudnika». - Almaty: TOO
«KazNIICvetmet», 2022. - 347 s. {[}in Russian{]}

6. Bekbergenov D.K., Zejnullin A.A. Razvitie tehnologii povtornoj
dobychi rudy iz mezhdukamernyh celikov, ostavlennyh v zonah obrushenija
na zavershajushhej stadii otrabotki Zhezkazganskogo mestorozhdenija //
Gornyj informacionno-analiticheskij bjulleten'. - 2025. -
№ 2 (specvypusk 4). - 20 s. {[}in Russian{]}

7. Demin V., Kalinin A., Tomilova N., Tomilov A., Akpanbayeva A.,
Shokarev D., Popov A., Advanced digital modeling of stress--strain
behavior in rock masses to ensure stability of underground mine
workings//Civil Engineering Journal. - 2025.-Vol.11(3). -P.1072-1087.
\href{https://doi.org/10.28991/CEJ-2025-011-03-014}{DOI
10.28991/CEJ-2025-011-03-014}.

8. Ahanov T. M. Sozdanie vysokojeffektivnyh tehnologij razrabotki
mestorozhdenij v uslovijah podzemnyh rudnikov TOO «Korporacija Kazahmys»
// Gornyj zhurnal Kazahstana. - 2009. - № 11. - S.26--30. {[}in
Russian{]}

9. Gerasimenko V. I., Makarov A. B., Zoteev O. V. Analiz
geomehanicheskogo sostojanija vyrabotannyh prostranstv Zhezkazganskogo
mestorozhdenija. Chast'{} 1. - Karaganda - Zhezkazgan:
GGU TOO «Korporacija Kazahmys, 2010. - 52 s. {[}in Russian{]}

10. Nizamova G. Z., Gajfullina M. M. Korreljacionno-regressionnyj analiz
jeffektivnosti ispol' zovanija investicionnyh resursov
neftjanoj kompanii // Vestnik UGNTU. Nauka, obrazovanie, jekonomika.
Serija: Jekonomicheskie nauki. - 2021. - T.9(3). -C.27-38 DOI
10.26425/2309-3633-2021-9-3-27-38. {[}in Russian{]}

11. Samsonov A. V., Vasil' ev I. V. Analiz investicionnoj
privlekatel' nosti predprijatij neftegazovoj otrasli v
Rossii // Moskovskij jekonomicheskij zhurnal. -- 2020. - № 2. - S.
439-450. DOI 10.24411/2413-046x-2020-10088. {[}in Russian{]}

12. Hastie T., Tibshirani R., Friedman J. The Elements of Statistical
Learning: Data Mining, Inference, and Prediction. - 2nd ed. - New York:
Springer. - 2009. - 745 p.
\end{refs}

\begin{info}
\hspace{1em}\emph{{\bfseries Сведения об авторах}}

Бекбергенов Д.К.- кандидат технических наук, Институт горного дела им.
Д.А.Кунаева, Алматы, Казахстан, е-mail: kdbekbergen@mail.ru;

Савич И. Н.- доктор технических наук, профессор, НИТУ МИСИС, Москва,
Россия, е-mail: tpr\_msmu@mail.ru;

Зейнуллин А.А.- доктор технических наук, профессор, АО «Казахский
университет технологии и бизнеса им.К.Кулажанова», Астана, Казахстан,
е-mail: karim\_57@mail.ru;

Бимурат Ж.- PhD, Институт горного дела им. Д.А.Кунаева, Алматы,
Казахстан, е-mail: zhbimurat@igd,com,kz;

Жанакова Р.К.- PhD ассоциированный профессор, Казахский национальный
исследовательский технический университет им. К. И. Сатпаева, Алматы,
Казахстан, е-mail: zhanakova\_raisa@mail.ru.

\hspace{1em}\emph{{\bfseries Information about the author}}

Bekbergenov D.K. - candidate of technical sciences, D.A. Kunaev
Institute of Mining, Almaty, Kazakhstan, e-mail: kdbekbergen@mail.ru;

Savich I.N. - Doctor of Technical Sciences, Professor, NUST MISIS,
Moscow, Russia, e-mail: tpr\_msmu@mail.ru;

Zeinullin A.A. - Doctor of Technical Sciences, Professor, JSC « K.
Kulazhanov Kazakh University of Technology and Business », Astana,
Kazakhstan, e-mail: karim\_57@mail.ru;

Zhanar Bю - PhD, D.A. Kunaev Institute of Mining, Almaty, Kazakhstan,
e-mail: zhbimurat@igd,com,kz;

Zhanakova R.K.- PhD Assoc. professor, K. I. Satpayev Kazakh National
Research Technical University, Almaty, Kazakhstan, e-mail:
zhanakova\_raisa@mail.ru.
\end{info}

